%!TEX root = main.tex

\subsection{Motivation and Design Rationale}

Phase~2 subjects the GlobalFlow network to six scenarios drawn directly
from the assignment specification. While these scenarios cover a broad
range of disruption types : tariff shocks, energy crises, node loss,
and corridor blockades. They share a common geographic focus on the
East Asian hub cluster and transatlantic trade lanes.
A structural blind spot remains: the network's dependence on the two
great maritime chokepoints that underpin global seaborne trade, namely
the Suez and Panama canals.

Together, these two canals handle approximately 12\% of world trade
by volume \cite{unctad2024}. Disruptions to either have historically
propagated rapidly through global supply chains, raising freight rates,
extending transit times, and forcing carriers to reroute through
significantly longer alternatives. The 2021 Ever Given grounding
blocked the Suez Canal for six days, triggering an estimated \$9.6bn
per day in delayed trade \cite{everGiven2021}. The 2023--2024 Houthi
attack campaign on Red Sea shipping diverted more than 50\% of normal
Suez traffic onto the Cape of Good Hope route within three months
\cite{drewry2024}. Simultaneously, a severe drought reduced Panama
Canal water levels to their lowest in recorded history, forcing
draft restrictions that cut effective throughput by nearly 40\%
\cite{panama2024}.

These events are not hypothetical tail risks. They are recent,
documented, and, crucially for GlobalFlow, they affect different
parts of the network simultaneously. Suez disruptions primarily
penalise Europe--Asia and Europe--MiddleEast flows, while Panama
disruptions cut across Americas--Asia and Americas--Europe corridors.
The Phase~2 scenarios do not capture this geography.

\subsection{Scenario Design}
\label{sec:phase3_design}

We design six new scenarios organised into two groups of three,
following a \emph{decomposed stress-testing} logic: each group
isolates first the cost dimension, then the capacity dimension, then
their combination. This structure allows us to disentangle two
economically distinct mechanisms of canal disruption.

\paragraph{Arc identification methodology.}
The identification of affected arcs is performed dynamically at
runtime using the zone-pair geography of the network. Each arc in
the instance is annotated with a \texttt{zone\_from} and
\texttt{zone\_to} pair indicating the regions it connects. For the
Suez corridor, we identify all sea-mode arcs crossing the zone pairs
that transit the Red Sea, Mediterranean, and Indian Ocean passages:
Europe$\leftrightarrow$Asia, Europe$\leftrightarrow$MiddleEast,
MiddleEast$\leftrightarrow$Asia and EastEurope$\rightarrow$MiddleEast. This yields 57~sea arcs.

For Panama-corridor arcs are identified as all sea arcs connecting 
Americas$\leftrightarrow$Asia, SouthAmerica$\leftrightarrow$Asia, and 
Americas$\leftrightarrow$SouthAmerica 
zone pairs (65 arcs total). The last category is included under the 
assumption that intra-American maritime freight primarily routes via 
Pacific coast ports, making it Canal-dependent.
This yields 122~sea arcs. 

The two sets are verified to be mutually exclusive
at runtime, ensuring no double-counting. This approach grounds the
scenario design firmly in the actual network topology rather than
relying on hardcoded arc lists, making the analysis robust to future
instance changes.

\paragraph{Cost shock (variants .1).}
When a canal remains open but vessels must reroute, around the Cape
of Good Hope for Suez, or Cape Horn for Panama, variable transport
costs rise sharply due to the additional distance (approximately
9\,000\,km and 12\,000\,km respectively). Arc capacities are
unaffected: ships divert, they do not disappear. This variant models
the carrier response to insecurity without assuming physical closure.
The varibale cost of all affected sea arcs increases by $80\%$ ($\times 1.80$). 

\paragraph{Capacity reduction (variants .2).}
Alternatively, a canal may remain legally open but physically
constrained, by military inspection regimes, draft restrictions due
to drought, or insurance-driven self-exclusion by major carriers.
In this case, the throughput of affected arcs falls while unit costs
for vessels that do transit remain near baseline. We model this as a
$60\%$ capacity reduction ($\times 0.40$) on all sea arcs transiting
the respective canal corridor.

\paragraph{Combined shock (variants .3).}
The most severe scenario applies both mechanisms simultaneously:
rerouting costs rise \emph{and} canal throughput is constrained.
This corresponds to the situation observed during the 2024 Red Sea
crisis, where carriers faced both elevated insurance premiums forcing
rerouting and a partial congestion of remaining Suez traffic due to
the residual vessels that continued transiting under escort.

\subsection{Scenario Descriptions}
\label{sec:phase3_scenarios}

Table~\ref{tab:canal_scenarios} summarises the six scenarios.

\begin{table}[H]
\centering
\caption{Phase~3 canal disruption scenarios.}
\label{tab:canal_scenarios}
\small
\begin{tabular}{llcc}
\toprule
\textbf{ID} & \textbf{Shock} & \textbf{Cost} & \textbf{Cap.} \\
\midrule
S4.1 & Suez cost shock       & $\times 1.80$ & $\times 1.00$ \\
S4.2 & Suez capacity cut     & $\times 1.00$ & $\times 0.40$ \\
S4.3 & Suez combined         & $\times 1.80$ & $\times 0.40$ \\
S5.1 & Panama cost shock     & $\times 1.80$ & $\times 1.00$ \\
S5.2 & Panama capacity cut   & $\times 1.00$ & $\times 0.40$ \\
S5.3 & Panama combined       & $\times 1.80$ & $\times 0.40$ \\
\bottomrule
\end{tabular}
\end{table}

\paragraph{S4.1 : Suez cost shock.}
Motivated by the 2024 Houthi attack campaign on Red Sea shipping
\cite{houthi2024}, this scenario raises the variable cost of all
76~sea arcs crossing the Suez corridor by $80\%$.
The affected zone pairs include Europe$\leftrightarrow$Asia,
Europe$\leftrightarrow$MiddleEast, MiddleEast$\leftrightarrow$Asia and EastEurope$\rightarrow$MiddleEast.
Carriers reroute around the Cape of Good Hope, adding roughly
9\,000\,km and 10--14 transit days. Arc capacities are unchanged:
all shipping capacity remains available, at higher cost.

\paragraph{S4.2 : Suez capacity reduction.}
Inspired by the Ever Given grounding (2021) and partial inspection
regimes, this scenario reduces the shared capacity of all 76~Suez
sea arcs by $60\%$ while leaving variable costs at baseline.
The $\times 0.40$ multiplier reflects a scenario where roughly
two-thirds of normal Suez traffic is diverted or blocked, while the
remaining third continues transiting at normal unit cost.

\paragraph{S4.3 : Suez combined shock.}
This scenario simultaneously applies the cost increase of~S4.1 and
the capacity restriction of~S4.2, representing the 2024 Red Sea
crisis at peak intensity: rerouting is both expensive (elevated
insurance, longer routes) and constrained (reduced throughput on
remaining Suez corridor arcs).

\paragraph{S5.1 : Panama cost shock.}
Motivated by the 2023--2024 drought that reduced canal water levels
to historic lows \cite{panama2024}, and by the geopolitical tensions
surrounding the canal's administration following US--Panama friction
in late 2024 \cite{panamaGeo2025}, this scenario raises variable
costs on all 122~Panama sea arcs by $80\%$.
Affected corridors include Americas$\leftrightarrow$Asia,
SouthAmerica$\leftrightarrow$Asia, and intra-Americas sea routes Americas$\leftrightarrow$SouthAmerica.
The $80\%$ premium reflects the additional distance and time cost of
rerouting via Cape Horn.

\paragraph{S5.2 : Panama capacity reduction.}
During the 2023–2024 drought the Panama Canal Authority implemented draft 
restrictions that lowered allowable vessel drafts and reduced daily transits, 
which significantly curtailed canal capacity; however, the relationship between
 draft limits and vessel deadweight varies by ship type and is not a uniform 40\%
  reduction in DWT. For modelling purposes, this scenario assumes a 60\% 
  reduction in effective capacity (×0.40) across all Panama-corridor sea arcs 
  to represent a case where drought or deliberate traffic management sharply 
  imits throughput while the canal remains nominally open.

\paragraph{S5.3 : Panama combined shock.}
The most severe Panama scenario combines the cost increase of~S5.1
with the capacity restriction of~S5.2: a drought-driven throughput
collapse paired with a punitive transit-fee hike, as debated by the
Panamanian government in late 2024 amid pressure to raise canal
revenues.

\subsection{Phase 3 Results}
\label{sec:phase3_results}

All six canal scenarios are structurally feasible: the $\times 0.40$
capacity reduction, while severe, does not exhaust the outgoing arc
capacity of any single supplier once the full network is considered.
This contrasts with the node-removal infeasibility of Scenarios S1
and~S3, where a hub is physically deleted. Here, canal arcs are
capacity-constrained but not removed, and the solver can always
route residual flow through alternative corridors, at higher cost.
No penalty slack is required and all demand is met under every
strategy.

Table~\ref{tab:phase3_results} reports the logistics cost achieved by
each strategy across the six scenarios, directly comparable to
$Z^* = \$4{,}454{,}801$.

\begin{table}[H]
\centering
\caption{Phase~3 logistics costs (\$). $\Delta Z = Z - Z^*$.
         Flex.\ value $= Z_R - Z_A$.}
\label{tab:phase3_results}
\resizebox{\linewidth}{!}{%
\footnotesize
\setlength{\tabcolsep}{5pt}
\begin{tabular}{@{}lccrrrrr@{}}
\toprule
\textbf{Sc.} & \textbf{Canal} & \textbf{Shock}
  & \multicolumn{1}{c}{$Z_R$}
  & \multicolumn{1}{c}{$Z_A$}
  & \multicolumn{1}{c}{$Z_F$}
  & \multicolumn{1}{c}{$\Delta Z_F$ (\%)}
  & \multicolumn{1}{c}{Flex.\ value} \\
\midrule
S4.1 & Suez   & C     & 4{,}764{,}699 & 4{,}764{,}699 & 4{,}764{,}699 & $+6.96\%$ & \$0 \\
S4.2 & Suez   & K     & 4{,}463{,}450 & 4{,}460{,}265 & 4{,}460{,}266 & $+0.12\%$ & \$3{,}185 \\
S4.3 & Suez   & C+K   & 4{,}798{,}860 & 4{,}798{,}860 & 4{,}798{,}860 & $+7.72\%$ & \$0 \\
S5.1 & Panama & C     & 4{,}682{,}184 & 4{,}674{,}317 & 4{,}673{,}500 & $+4.91\%$ & \$7{,}868 \\
S5.2 & Panama & K     & 4{,}465{,}620 & 4{,}462{,}902 & 4{,}462{,}851 & $+0.18\%$ & \$2{,}718 \\
S5.3 & Panama & C+K   & 4{,}682{,}184 & 4{,}674{,}317 & 4{,}673{,}500 & $+4.91\%$ & \$7{,}868 \\
\bottomrule
\end{tabular}%
}
\begin{minipage}{\linewidth}
\vspace{3pt}
\footnotesize
Shock codes: C\,=\,cost only, K\,=\,capacity only, C+K\,=\,combined.
All demand is met under every strategy and scenario.
\end{minipage}
\end{table}

\paragraph{Monotonicity check.}
The expected ordering $Z_F \leq Z_A \leq Z_R$ holds in all six
scenarios. For S4.1 and S4.3, all three strategies tie exactly
($Z_R = Z_A = Z_F$): the existing warehouse footprint is already
optimal under a pure Suez cost shock, and no reconfiguration
improves upon it. In the Panama group, F retains a small edge over
A (\$51--\$817) by deactivating one baseline arc whose sunk cost
is no longer justified.

\subsubsection*{Suez group (S4.1--S4.3)}

\paragraph{S4.1 : Suez cost shock.}
The $+80\%$ surcharge on the 39~affected Suez sea arcs raises
$Z_F$ to \$4{,}764{,}699 ($+6.96\%$). All three strategies tie:
the baseline network absorbs the cost shock without any
reconfiguration, and the flexibility value is exactly \$0.
A rerouting of all Suez flows around the Cape of Good Hope is optimal, 
but the baseline solution already routes all affected flow
only through Asia$\leftrightarrow$Europe legs.
\begin{table}[H]
\centering
\caption{Cost decomposition --- scenario S4.1 (\$).}
\label{tab:cost_decomp_s4_1}
\resizebox{\linewidth}{!}{%
\begin{tabular}{lrr}
\toprule
\textbf{Component} & \textbf{Baseline} & \textbf{R, A, F} \\
\midrule
\multicolumn{3}{l}{\textit{Fixed costs}} \\
~~WH opening (sunk) & \$22,046 & \$22,046 \\
~~WH opening (new)  & \$0 & \$0 \\
~~Arc activation (sunk) & \$19,638 & \$19,638 \\
~~Arc activation (new)  & \$0 & \$0 \\
\midrule
\multicolumn{3}{l}{\textit{Variable transport costs by product}} \\
~~A (Fertilizers) & \$638,198 & \$639,128 \\
~~B (Semiconductors) & \$2,814,741 & \$3,009,410 \\
~~C (Battery Comp.) & \$960,176 & \$1,074,477 \\
\midrule
\textbf{Total $Z$} & \$4,454,801 & \$4,764,699 \\
\midrule
$\Delta Z$ vs $Z^*$ & --- & \$+309,899 \\
$\Delta Z$ (\%) & --- & +6.96\% \\
Flex.\ value $Z_R-Z_A$ & --- & \$0 \\
\bottomrule
\end{tabular}
}
\end{table}

\paragraph{S4.2 : Suez capacity reduction.}
The $\times 0.40$ capacity cut is almost entirely absorbed:
$Z_F = \$4{,}460{,}266$ ($+0.12\%$), with all demand met.
Unlike the previous model version, the residual corridor capacity
is sufficient to route all supplier volumes, meaning no supplier
is bottlenecked at origin. The small flexibility value
(\$3{,}185) reflects marginal flow rebalancing at the warehouse
layer. This is the Suez scenario GlobalFlow is best positioned
to survive.
\begin{table}[H]
\centering
\caption{Cost decomposition --- scenario S4.2 (\$).}
\label{tab:cost_decomp_s4_2}
\resizebox{\linewidth}{!}{%
\begin{tabular}{lrrrr}
\toprule
\textbf{Component} & \textbf{Baseline} & \textbf{R} & \textbf{A} & \textbf{F} \\
\midrule
\multicolumn{5}{l}{\textit{Fixed costs}} \\
~~WH opening (sunk) & \$22,046 & \$22,046 & \$22,046 & \$0 \\
~~WH opening (new)  & \$0 & \$0 & \$0 & \$22,046 \\
~~Arc activation (sunk) & \$19,638 & \$19,638 & \$19,638 & \$0 \\
~~Arc activation (new)  & \$0 & \$0 & \$1,671 & \$21,309 \\
\midrule
\multicolumn{5}{l}{\textit{Variable transport costs by product}} \\
~~A (Fertilizers) & \$638,198 & \$638,198 & \$638,198 & \$638,198 \\
~~B (Semiconductors) & \$2,814,741 & \$2,823,198 & \$2,818,343 & \$2,818,343 \\
~~C (Battery Comp.) & \$960,176 & \$960,368 & \$960,368 & \$960,369 \\
\midrule
\textbf{Total $Z$} & \$4,454,801 & \$4,463,450 & \$4,460,265 & \$4,460,266 \\
\midrule
$\Delta Z$ vs $Z^*$ & --- & \$+8,649 & \$+5,465 & \$+5,465 \\
$\Delta Z$ (\%) & --- & +0.19\% & +0.12\% & +0.12\% \\
Flex.\ value $Z_R-Z_A$ & --- & \multicolumn{2}{c}{\$3,185} & --- \\
\bottomrule
\end{tabular}
}
\end{table}

\paragraph{S4.3 : Suez combined shock.}
Combining the $+80\%$ cost premium with the $\times 0.40$
capacity cut yields $Z_F = \$4{,}798{,}860$ ($+7.72\%$), with all
three strategies again tying exactly. The capacity constraint
does not add binding pressure beyond the cost shock alone:
once the solver reroutes away from the most expensive Suez arcs,
the residual capacity on those arcs is no longer the active
constraint. The combined disruption cost ($+7.72\%$) is
slightly larger than S4.1 ($+6.96\%$) but not additive, 
the interaction between the two mechanisms is sublinear.

\begin{table}[H]
\centering
\caption{Cost decomposition --- scenario S4.3 (\$). All three strategies (R, A, F) yield identical costs.}
\label{tab:cost_decomp_s4_3}
\resizebox{\linewidth}{!}{%
\begin{tabular}{lrr}
\toprule
\textbf{Component} & \textbf{Baseline} & \textbf{R / A / F} \\
\midrule
\multicolumn{3}{l}{\textit{Fixed costs}} \\
~~WH opening & \$22,046 & \$22,046 \\
~~Arc activation & \$19,638 & \$19,638 \\
\midrule
\multicolumn{3}{l}{\textit{Variable transport costs by product}} \\
~~A (Fertilizers) & \$638,198 & \$639,128 \\
~~B (Semiconductors) & \$2,814,741 & \$3,023,541 \\
~~C (Battery Comp.) & \$960,176 & \$1,094,505 \\
\midrule
\textbf{Total $Z$} & \$4,454,801 & \$4,798,860 \\
\midrule
$\Delta Z$ vs $Z^*$ & --- & \$+344,059 \\
$\Delta Z$ (\%) & --- & +7.72\% \\
\bottomrule
\end{tabular}
}
\end{table}

\subsubsection*{Panama group (S5.1--S5.3)}

\paragraph{S5.1 : Panama cost shock.}
The $+80\%$ cost shock on the 65~Panama-corridor arcs produces
$Z_F = \$4{,}673{,}500$ ($+4.91\%$), materially less than the
Suez equivalent ($+6.96\%$). The flexibility value \$7{,}868 is
the largest in the cost-shock group: adaptation activates two
bypass arcs (A223 W14$\rightarrow$W12 sea; A224
W12$\rightarrow$C1 road) that route Americas-bound flows through
non-Panama lanes. The Panama corridor carries a smaller baseline
share than Suez, and natural detours are available.
\begin{table}[H]
\centering
\caption{Cost decomposition --- scenario S5.1 (\$).}
\label{tab:cost_decomp_s5_1}
\resizebox{\linewidth}{!}{%
\begin{tabular}{lrrrr}
\toprule
\textbf{Component} & \textbf{Baseline} & \textbf{R} & \textbf{A} & \textbf{F} \\
\midrule
\multicolumn{5}{l}{\textit{Fixed costs}} \\
~~WH opening (sunk) & \$22,046 & \$22,046 & \$22,046 & \$0 \\
~~WH opening (new)  & \$0 & \$0 & \$0 & \$22,046 \\
~~Arc activation (sunk) & \$19,638 & \$19,638 & \$19,638 & \$0 \\
~~Arc activation (new)  & \$0 & \$0 & \$2,461 & \$21,282 \\
\midrule
\multicolumn{5}{l}{\textit{Variable transport costs by product}} \\
~~A (Fertilizers) & \$638,198 & \$638,566 & \$637,799 & \$637,799 \\
~~B (Semiconductors) & \$2,814,741 & \$2,885,746 & \$2,883,843 & \$2,883,843 \\
~~C (Battery Comp.) & \$960,176 & \$1,116,188 & \$1,108,530 & \$1,108,530 \\
\midrule
\textbf{Total $Z$} & \$4,454,801 & \$4,682,184 & \$4,674,317 & \$4,673,500 \\
\midrule
$\Delta Z$ vs $Z^*$ & --- & \$+227,384 & \$+219,516 & \$+218,700 \\
$\Delta Z$ (\%) & --- & +5.10\% & +4.93\% & +4.91\% \\
Flex.\ value $Z_R-Z_A$ & --- & \multicolumn{2}{c}{\$7,868} & --- \\
\bottomrule
\end{tabular}
}
\end{table}

\paragraph{S5.2 : Panama capacity reduction.}
The $\times 0.40$ capacity cut on Panama arcs is almost
entirely absorbed: $Z_F = \$4{,}462{,}851$ ($+0.18\%$). The
baseline Panama flows are well below the residual corridor
capacity, so the constraint never binds. The flexibility value
(\$2{,}718) is correspondingly small. This is the canal scenario
GlobalFlow is best positioned to absorb.
\begin{table}[H]
\centering
\caption{Cost decomposition --- scenario S5.2 (\$).}
\label{tab:cost_decomp_s5_2}
\resizebox{\linewidth}{!}{%
\begin{tabular}{lrrrr}
\toprule
\textbf{Component} & \textbf{Baseline} & \textbf{R} & \textbf{A} & \textbf{F} \\
\midrule
\multicolumn{5}{l}{\textit{Fixed costs}} \\
~~WH opening (sunk) & \$22,046 & \$22,046 & \$22,046 & \$0 \\
~~WH opening (new)  & \$0 & \$0 & \$0 & \$22,046 \\
~~Arc activation (sunk) & \$19,638 & \$19,638 & \$19,638 & \$0 \\
~~Arc activation (new)  & \$0 & \$0 & \$1,718 & \$20,539 \\
\midrule
\multicolumn{5}{l}{\textit{Variable transport costs by product}} \\
~~A (Fertilizers) & \$638,198 & \$638,566 & \$637,799 & \$637,799 \\
~~B (Semiconductors) & \$2,814,741 & \$2,885,746 & \$2,883,843 & \$2,883,843 \\
~~C (Battery Comp.) & \$960,176 & \$1,116,188 & \$1,108,530 & \$1,108,530 \\
\midrule
\textbf{Total $Z$} & \$4,454,801 & \$4,465,620 & \$4,462,902 & \$4,462,851 \\
\midrule
$\Delta Z$ vs $Z^*$ & --- & \$+10,820 & \$+8,101 & \$+8,050 \\
$\Delta Z$ (\%) & --- & +0.24\% & +0.18\% & +0.18\% \\
Flex.\ value $Z_R-Z_A$ & --- & \multicolumn{2}{c}{\$2,718} & --- \\
\bottomrule
\end{tabular}
}
\end{table}


\paragraph{S5.3 : Panama combined shock.}
S5.3 produces identical costs to S5.1 ($Z_F = \$4{,}673{,}500$,
$\Delta Z_F = +4.91\%$, flex.\ value \$7{,}868). Once the solver
reroutes away from Panama in response to the $+80\%$ cost shock,
the $\times 0.40$ capacity cap on the corridor becomes
non-binding: the optimal solution already avoids Panama, so
restricting its capacity costs nothing additional. The Panama
vulnerability is a \emph{price} vulnerability, not a
\emph{throughput} vulnerability.
\begin{table}[H]
\centering
\caption{Cost decomposition --- scenario S5.3 (\$).}
\label{tab:cost_decomp_s5_3}
\resizebox{\linewidth}{!}{%
\begin{tabular}{lrrrr}
\toprule
\textbf{Component} & \textbf{Baseline} & \textbf{R} & \textbf{A} & \textbf{F} \\
\midrule
\multicolumn{5}{l}{\textit{Fixed costs}} \\
~~WH opening (sunk) & \$22,046 & \$22,046 & \$22,046 & \$0 \\
~~WH opening (new)  & \$0 & \$0 & \$0 & \$22,046 \\
~~Arc activation (sunk) & \$19,638 & \$19,638 & \$19,638 & \$0 \\
~~Arc activation (new)  & \$0 & \$0 & \$2,461 & \$21,282 \\
\midrule
\multicolumn{5}{l}{\textit{Variable transport costs by product}} \\
~~A (Fertilizers) & \$638,198 & \$638,566 & \$637,799 & \$637,799 \\
~~B (Semiconductors) & \$2,814,741 & \$2,885,746 & \$2,883,843 & \$2,883,843 \\
~~C (Battery Comp.) & \$960,176 & \$1,116,188 & \$1,108,530 & \$1,108,530 \\
\midrule
\textbf{Total $Z$} & \$4,454,801 & \$4,682,184 & \$4,674,317 & \$4,673,500 \\
\midrule
$\Delta Z$ vs $Z^*$ & --- & \$+227,384 & \$+219,516 & \$+218,700 \\
$\Delta Z$ (\%) & --- & +5.10\% & +4.93\% & +4.91\% \\
Flex.\ value $Z_R-Z_A$ & --- & \multicolumn{2}{c}{\$7,868} & --- \\
\bottomrule
\end{tabular}
}
\end{table}

\paragraph{Suez versus Panama.}
The cross-canal comparison is unambiguous on the cost dimension.
The Suez cost shock (S4.1, $+6.96\%$) is $1.4\times$ more damaging
than its Panama counterpart (S5.1, $+4.91\%$), reflecting the larger
share of GlobalFlow's baseline flow that transits the Red Sea
corridor. On the capacity dimension, the gap narrows sharply: both
S4.2 ($+0.12\%$) and S5.2 ($+0.18\%$) are absorbed near-costlessly,
confirming that the $\times 0.40$ residual capacity is sufficient
for all supplier volumes in the updated instance. The key asymmetry
lies in the \emph{combined} shock: S4.3 ($+7.72\%$) exceeds S5.3
($+4.91\%$) by a wider margin than the cost-only comparison alone,
because the Suez corridor supports a larger volume of
cost-sensitive flow that gets squeezed by the joint constraint.
Suez is the more exposed chokepoint; Panama is the more
\emph{resilient} one.

\subsubsection*{Cross-phase comparison}

\paragraph{Magnitude of canal shocks versus Phase~2 shocks.}
On comparable $\Delta Z_F$ percentages, the canal cost shocks
sit between T1 and T2: S4.1 ($+6.96\%$) falls between T1
($+3.49\%$) and T2 ($+10.64\%$); S5.1 ($+4.91\%$) is close to T1.
The capacity-only shocks (S4.2 $+0.12\%$, S5.2 $+0.18\%$) are
the mildest disruptions in the entire scenario library --- milder
even than T1. The topological shocks of Phase~2 (S1 $+81.6\%$,
S3 $+89.6\%$) remain an order of magnitude more damaging than
any canal scenario: removing a hub severs supplier connectivity
in a way that no capacity reduction on a sea corridor can replicate,
since alternative arcs remain available in all Phase~3 scenarios.

\paragraph{Where adaptation pays.}
The flexibility value $Z_R - Z_A$ tells a different story than
in Phase~2. For the Suez cost shocks (S4.1, S4.3), the flexibility
value is exactly \$0 --- the baseline network is already optimally
positioned and reconfiguration adds nothing. For the Panama cost
scenarios (S5.1, S5.3), adaptation yields \$7{,}868 by activating
two bypass arcs, a modest but non-trivial gain. The capacity-only
scenarios (S4.2 \$3{,}185; S5.2 \$2{,}718) sit in between.
Across the entire Phase~3 library, flexibility values are small
relative to T3 (\$98{,}244), confirming that canal disruptions ---
even severe ones --- do not fundamentally challenge the baseline
network topology the way a hub-handling shock does.

\subsection{Conclusion}
\label{sec:phase3_conclusion}

The Phase~3 stress test yields three headline findings. First,
all six canal scenarios are structurally feasible: the $\times 0.40$
capacity reduction leaves sufficient corridor throughput for all
supplier volumes, and every demand unit is met under every strategy.
This is a materially more optimistic picture than the previous model
version and reflects a corrected arc identification (39 Suez arcs,
65 Panama arcs) grounded in the actual instance topology.

Second, the Suez corridor is the more exposed of the two chokepoints.
Its cost-only shock ($+6.96\%$) is $1.4\times$ more damaging than
Panama's ($+4.91\%$), and its combined shock ($+7.72\%$) widens
the gap further. However, both canals remain manageable disruptions
relative to Phase~2 topological shocks: no canal scenario exceeds
$+8\%$, whereas H3 closure costs $+81.6\%$.

Third, adaptation (strategy~A) adds value only selectively.
For Suez cost shocks, the flexibility value is exactly zero ---
the baseline routing is already optimal. For Panama cost shocks,
two bypass arcs (A223, A224) unlock \$7{,}868 in savings. The
strategic implication is that the GlobalFlow network is
\emph{well-positioned against canal disruptions as currently
configured}, but that the handful of inter-warehouse arcs that
adaptation keeps activating across scenarios --- A199, A223,
A224 --- remain strong candidates for pre-activation in the
baseline at negligible fixed cost.

This conclusion reinforces the Phase~2 finding: the network's
primary vulnerability is not geographic breadth but topological
concentration. The H3 hub and the two suppliers that route
exclusively through it (S6 Yokohama, S9 Chengdu) are exposed
to Phase~2 and Phase~3 shocks alike. A genuinely robust
GlobalFlow network addresses this concentration first;
canal resilience, by comparison, is already largely built in.